\section{多目标优化：前沿形状与初始化}
\subsection{先定义可行解和收入，才比较算法}
基础NSGA单独运行20司机/20合成OD，10×10 km欧氏空间；两个目标为最小总接驾距离与最大总送客km，不引入进阶车费。所有边可行，匹配可选；每个固定$\lambda$按$U=\sum d_{OD}-\lambda\sum d_{pickup}$选解，与完全相同可行图的精确最大效用KM比较。五seed结果及前沿见图\ref{fig:basicfront}。该静态效用实例近似比为$U_{OPT}/U_{EA}$，不同于上一节在线延迟比；无标量偏好的一整条Pareto前沿没有唯一竞争比。
\datafigure{evolution_basic_scalar_reference}{基础部分独立的欧氏OD实验，五个进化seed，100种群/300代。车费目标仅为送客km；标量化参考来自同一可行图的精确求解，不是在线理论界。}{fig:basicfront}

进阶整数编码长度为订单数，第$j$位0表示不派，其余值指向该单最近K个可行司机；按请求顺序解码，司机冲突时轮转候选，耗尽则不派。所有算法共享这个半径/候选图与单次容量。统一最小化目标$F=(-\sum fare,\sum pickup)$，第三目标加入$-\#served$，因此空匹配是合法端点，不用伪造服务数提高收入。

\subsection{非支配保留和共同评价标尺}
NSGA-II以非支配层、逐目标归一化拥挤距离和精英保留筛选\cite{nsga2}；NSGA-III保留非支配层，把临界层截断换成理想点/超平面归一化、参考方向关联和小生境保留\cite{nsga3}。这些机制按教师整数编码作教学实现。二目标100种群/100方向；三目标按Notebook的105种群/105方向，奇数种群必须生成实际105个子代而非104个。交叉概率0.9，每位变异概率$1/m$；100种群/300代的实现实际评价30,100个个体（包括初代）。

HV使用所有方法共同的预先界：收入上界由同可行图精确最大收入匹配求得，接驾上界为容量×3 km，第三目标服务上界为容量；各维归一化后参考点统一1.1。HV可大于1，因为参考盒边长1.1，而不是把它解释为百分比。Spacing采用不同规范非支配点之间最近L1距离的变异系数；不足两点或均距0时无定义。

独立审核发现，加权KM的离散权重输出可能包含支配点，以及求和次序带来的约$10^{-13}$差异。保留原始输出，报告端对所有方法统一按10位小数规范去重后取非支配集；8/10/12位敏感性与另写box-union核验确认其影响范围。不能只对某一个算法换标尺使其看起来更好。

\datafigure{evolution_front_and_convergence}{18点inner训练批次，5个进化seed；共同可行边与固定HV标尺。收敛阴影为5seed最小--最大范围，KM只求离散权重方向，不保证完整前沿。}{fig:frontconv}
\begin{table}[htbp]\centering\small
\caption{18点相同开发批次；5seed均值±样本标准差，DTLZ2独立使用92/400设置。}\label{tab:evobaseline}
\begin{tabularx}{\linewidth}{Ylrrr}\toprule
问题 & 算法 & HV & Spacing & EA秒\\\midrule
派单2D & NSGA-II & 1.090244 $\pm$ 0.004460 & 0.505 & 2.455\\
派单2D & NSGA-III & 1.077092 $\pm$ 0.009530 & 0.938 & 3.601\\
派单3D & NSGA-II & 1.014845 $\pm$ 0.010260 & 0.621 & 2.560\\
派单3D & NSGA-III & 1.012638 $\pm$ 0.005689 & 0.688 & 4.220\\
DTLZ2 & NSGA-II & 0.707111 $\pm$ 0.006190 & 0.480 & 1.087\\
DTLZ2 & NSGA-III & 0.744325 $\pm$ 0.000235 & 0.297 & 1.841\\
\bottomrule\end{tabularx}
\tabnote{各问题使用各自预先定义的共同标尺，HV不是百分比，也不跨不同问题比高低。}
\end{table}


\subsection{加权和何时遗漏，三目标何时需要参考方向}
在该离散派单问题中，加权和KM能高效找到很多supported解，尤其补足高收入端；但非凸部分的unsupported解不可能成为某个线性权重最优。一个已知答案的最小化例子是$(0,2),(1,1.1),(2,0)$：中点不被支配，却无法胜过两端的加权和最优。这里还有限权重网格的覆盖问题，即便supported解也可能漏掉，不能把两种原因混为一谈。

三目标派单中的服务数与收入高度相关，前沿可能仍很窄；目标维数增加不保证NSGA-III胜出。DTLZ2提供真正铺开的球面几何，图\ref{fig:three}与图\ref{fig:dtlz}并置说明参考方向的适用差异。DTLZ2使用教师给定的离散编码实现及独立公式参考，不把它冒称任意连续全局优化保证。Spacing是局部分布离散度，需与HV、端点、误差一起判断。

\datafigure{evolution_three_objective_geometry}{实际三目标派单前沿：负号转换回正收入/服务数展示，三维和二维投影互为补充。第三维是真实服务数。}{fig:three}
DTLZ2五seed的平均球面径向误差$|\lVert F\rVert_2-1|$分别为II $0.0076465\pm0.0013867$、III $0.00007084\pm0.00005153$（样本标准差）。这项有已知球面参考的误差与Spacing提供互补信息。
\datafigure{evolution_dtlz2_geometry}{DTLZ2的seed0实测点与解析球面网格；多次运行的径向误差见正文，分布统计见表。5seed统计与派单场景独立，不能直接比较不同实例的绝对HV大小。}{fig:dtlz}

\subsection{先增加预算，再测试KM种子}
基础18点结果的高收入端不足，先比较种群200/代数300，以及种群100/代数600，二者约60,000次评价。随后保持原100/300预算，插入最多10个真实编码的KM解，其余仍为原初始化；验证$decode(encode(KM))$与原匹配完全一致。初始化额外41次标量求解、41次目标求和和耗时单列，没有藏在EA计时之外。

观察到KM初始化稳定改善端点和前沿后，才组合KM初始化与种群200，保留原基线、种群单项、KM单项和组合四个父项/移除对照（图\ref{fig:budget}）。没有同时叠加所有增强。最终留出只确认冻结的NSGA-II原100/300与同预算KM初始化，在未来日期池5个独立批次上成对，未据测试选择种群和权重。

\datafigure{evolution_budget_and_initialization}{基线、种群、代数、KM单项及有限组合的HV与总耗时；目标评价次数另见表。初始化成本纳入总耗时，重复seed不挑最好一次；时间为并发现场观测，不据此宣称严格性能加速。}{fig:budget}
\begin{table}[htbp]\centering\small
\caption{18点开发输入的单项、父项和组合，每项5个进化seed。}\label{tab:evobudget}
\begin{tabularx}{\linewidth}{lYrrrr}\toprule
算法 & 变化 & HV均值 & Spacing & 目标评价 & 总秒\\\midrule
II & 基线 & 1.090244 & 0.505 & 30,100 & 2.455\\
III & 基线 & 1.077092 & 0.938 & 30,100 & 3.601\\
II & 种群200 & 1.101432 & 0.617 & 60,200 & 4.294\\
III & 种群200 & 1.085327 & 1.180 & 60,200 & 5.651\\
II & 代数600 & 1.094778 & 0.501 & 60,100 & 3.296\\
III & 代数600 & 1.090552 & 0.931 & 60,100 & 4.504\\
II & KM初始 & 1.111427 & 0.470 & 30,100(+41) & 2.032\\
III & KM初始 & 1.111494 & 0.901 & 30,100(+41) & 3.021\\
II & KM+种群 & 1.112912 & 0.626 & 60,200(+41) & 4.519\\
III & KM+种群 & 1.112072 & 1.229 & 60,200(+41) & 5.939\\
\bottomrule\end{tabularx}
\tabnote{括号为额外KM初始化目标求和次数，另有相同次数标量求解；共同归一化成本分开记录。现场并发耗时不作严格微基准。}
\end{table}

\begin{table}[htbp]\centering\small
\caption{冻结后的5个未来日期池独立采样输入，算法预算相同。}\label{tab:evofuture}
\begin{tabularx}{\linewidth}{Yrrr}\toprule
输入seed & Teacher-II HV & KM初始化HV & 成对差\\\midrule
9010 & 1.042422 & 1.074096 & +0.031674\\
9011 & 1.052557 & 1.092004 & +0.039447\\
9012 & 1.067154 & 1.089801 & +0.022647\\
9013 & 1.059461 & 1.087311 & +0.027849\\
9014 & 1.072543 & 1.108278 & +0.035735\\
\bottomrule\end{tabularx}
\tabnote{未来输入seed9010--9014、进化seed10--14；每对在自身共同标尺下相减，不重新选择参数。}
\end{table}


\subsection{时段与最终可用折中}
8、13、18、23点从各自inner代理池抽批次，司机先验也限训练签到。同样司机数并不意味着同样可行订单数，前沿长度和高度需结合候选图解释（图\ref{fig:hourly}）。
\datafigure{evolution_hourly_fronts}{4时段、共同40司机/80订单与3km/K6规则；每时段保留5seed。来源时段、供需可行图与目标范围不同，不把前沿高度单独当算法优劣。}{fig:hourly}

膝点取归一化前沿与两端连线的最大垂距，不需平台先给硬阈值，但依赖目标尺度和端点，局部前沿变化会使推荐不稳定。约束法则在接驾预算下取最高收入，含义更容易说明，但预算需要真实业务依据。本次预先使用60 km批次总接驾预算作教学比较；在图示批次它不构成活跃约束，因而选最高收入EA解，不能称为已验证的实际SLA。两种规则及各seed实际选择点保留在结果中，不临时加权成一个所谓综合质量分。
